% s-124-p124-m003; original PDF 124
\Aphorism{s-124-p124-m003}{161}
Die alten Helden warben mit Siegen und Beuten um das Weib, wir mit Liedern und schönen Gefühlen; kein Wunder, dass wir den Gipfel des Glückes bei Weibern nicht mehr erreichen. Wir sprechen von „Waffen des Geistes“, von „Heroismus der Gesinnung“, ja wir nennen uns mit schrecklicher Naivität „Helden der Feder“, — der Feder, die eine verkleinerte Parodie des Speeres ist. In jedem rechten Weibe ist ein Letztes, das sich nur dem starken Manne ergiebt, beileibe nicht dem „geistreichen“: etwas von jenem Carmen-Idealismus, der die Männer in Gefahr und Tod schickt, um sie in ihrer Schönheit zu sehen. Und warum geht der Mann in Tod und Gefahr? Um das Weib in seiner Schönheit zu sehen, das Weib, das einem Sieger dankt. Wofür in aller Welt hätte es heute zu danken? Für ein Bändchen Lyrik, ein bestandenes Examen, einen ergriffenen Beruf — für eine Dummheit mehr an dem ohnehin schon so dummen \Foreign{masculinum}.
\par

